% !TEX root = ../main.tex

\section{Oscylator harmoniczny}
Jednym z podstawowych układów zarówno klasycznych, jak i kwantowych, jest oscylator harmoniczny. Opisany jest on hamiltonianem:
\begin{equation*}
	\ham = \inv{2m} \vu{p}^2 + \inv{2} m \omega^2 \vu{q}^2,
\end{equation*}
gdzie \(\vu{p}\) jest operatorem pędu, który w reprezentacji położeniowej ma postać \(\vu{p} = - i \hbar \pdv{}{q}\); \(\vu{q}\) jest operatorem położenia, \(m\) jest masą cząstki, a \(\omega\) jest częstotliwości drgań własnych i jest zależna od krzywizny potencjału, w którym znajduje się ta cząstka. W celu uproszczenia analizy zagadnienia wprowadzamy zmienne naturalne:
\begin{align*}
	\vu{Q} &= \sqrt{\frac{m\omega}{\hbar}} \vu{q}, &
	\vu{Q}^2 &= \frac{m\omega}{\hbar} \vu{q}^2, \\
	\pdv{}{Q} &= \sqrt{\frac{\hbar}{m\omega}} \pdv{}{q}, &
	\pdv[order=2]{}{Q} &= \frac{\hbar}{m\omega} \pdv[order=2]{}{q}, \\
	\vu{P} &= - i \pdv{}{Q} = \sqrt{\frac{\hbar}{m\omega}} \pdv{}{q}, &
	\vu{P}^2 &= - \pdv[order=2]{}{Q} = - \frac{\hbar}{m\omega} \pdv[order=2]{}{q}.
\end{align*}
Hamiltonian w tych zmiennych ma formę:
\begin{equation*}
	\ham = \frac{\hbar\omega}{2}\pab{\vu{Q}^2 + \vu{P}^2}.
\end{equation*}
Podobnie jak położenia i pędy spełniają zależność komutacyjną
\begin{equation*}
	\bab{\vu{q},\vu{p}} = i\hbar,
\end{equation*}
dla nowych zmiennych mamy:
\begin{equation*}
	\bab{\vu{Q},\vu{P}} = i.
\end{equation*}
Komutator operatorów o przeciwnych znakach:
\begin{equation*}
	\bab{\vu{Q} \pm i\vu{P}, \vu{Q} \mp i\vu{P}} = \pm 2,
\end{equation*}
co zostanie później wykorzystane do normalizacji operatorów.
Komutator hamiltonianu z nimi jest równy:
\begin{align*}
	\bab{\ham,\vu{Q}} = - i\hbar \vu{P}, \\
	\bab{\ham,\vu{P}} =   i\hbar \vu{Q},
\end{align*}
lub w pojedynczym równaniu:
\begin{equation*}
	\bab{\ham,\vu{Q} \pm i\vu{P}} = \mp \hbar\omega \pab{\vu{Q} \pm i\vu{P}}.
\end{equation*}

Operatory, które spełniają podobne zależności komutacyjne maja specyficzną własność --- są to operatory drabinkowe. Możemy znaleźć energię stanu \(\pab{\vu{Q} \pm i\vu{P}}\ket{n}\):
\begin{multline*}
	\ham \pab{\vu{Q} \pm i\vu{P}} \ket{n} =
	\pab{\pab{\vu{Q} \pm i\vu{P}} \ham + \bab{\ham,\vu{Q} \pm i\vu{P}}} \ket{n} = \\ =
	\pab{\vu{Q} \pm i\vu{P}} \ham \ket{n} \mp \hbar\omega \pab{\vu{Q} \pm i\vu{P}} \ket{n} =
	\pab{E_n \mp \hbar\omega} \pab{\vu{Q} \pm i\vu{P}} \ket{n}.
\end{multline*}
Rozpatrywany stan ma energię innego stanu, więc:
\begin{equation*}
	\pab{\vu{Q} \pm i\vu{P}}\ket{n} = C_n^{\pm} \ket{n \mp 1},
\end{equation*}
daje nam nowy stan o energii wyższej lub niższej o \(\hbar\omega\).

Definiujemy nowe operatory anihilacji i kreacji:
\begin{align*}
	\vu{a} &= \inv{\sqrt{2}} \pab{\vu{Q} + i\vu{P}}, &
	\vu{a}^{\dagger} &= \inv{\sqrt{2}} \pab{\vu{Q} - i\vu{P}}.
\end{align*}
Dodatkowo można znaleźć, że:
\begin{align*}
	\vu{a}^{\dagger}\vu{a} &= \inv{\hbar\omega}\ham - \inv{2}, &
	\ham &= \hbar\omega \pab{\vu{a}^{\dagger}\vu{a} + \inv{2}} =
	\hbar\omega \pab{\vu{a}\vu{a}^{\dagger} - \inv{2}}.
\end{align*}
Możemy zapisać poprzednio znalezione wzory wykorzystujac te operatory:
\begin{gather}\label{eq:osc_harm_wlasnosci_a}
	\bab{\vu{a},\vu{a}^{\dagger}} = 1, \\
	\bab{\ham,\vu{a}} = - \hbar\omega \vu{a}, \\
	\bab{\ham,\vu{a}^{\dagger}} = \hbar\omega \vu{a}^{\dagger}, \\
	\ham \vu{a} \ket{n} = \pab{E_n - \hbar\omega} \ket{n}, \\
	\ham \vu{a}^{\dagger} \ket{n} = \pab{E_n + \hbar\omega} \vu{a}^{\dagger} \ket{n}.
\end{gather}

Oscylator harmoniczny ma stan podstawowy:
\begin{equation*}
	\vu{a} \ket{0} = 0.
\end{equation*}
Energia tego stanu jest równa:
\begin{equation*}
	\ham \ket{0} =
	\hbar\omega \pab{\vu{a}^{\dagger}\vu{a} + \inv{2}} \ket{0} =
	\hbar\omega \vu{a}^{\dagger}\vu{a} \ket{0} + \inv{2}\hbar\omega \ket{0} =
	E_0 \ket{0},
\end{equation*}
stąd energia stanu podstawowego:
\begin{equation*}
	E_0 = \inv{2}\hbar\omega \ket{0}.
\end{equation*}
Korzystając z \eqref{eq:osc_harm_wlasnosci_a} możemy znaleźć energie stanów wyższych iteracyjnie co daje:
\begin{equation}\label{eq:osc_harm_En}
	E_n = \hbar\omega \pab{n + \inv{2}}.
\end{equation}

Znając energie możemy znaleźć stałe normalizacyjne sąsiednich stanów:
\begin{gather*}
	\vu{a} \ket{n} = C_n \ket{n-1}, \\
	\vab{C_n}^2 = \bra{n} \vu{a}^{\dagger}\vu{a} \ket{n} =
	\bra{n} \inv{\hbar\omega}\ham - \inv{2} \ket{n} =
	\inv{\hbar\omega} \hbar\omega \pab{n + \inv{2}} - \inv{2} = n, \\
	\vu{a}^{dagger} \ket{n} = D_n \ket{n-1}, \\
	\vab{D_n}^2 = \ket{n} \vu{a}\vu{a}^{\dagger} \ket{n} =
	\bra{n} \inv{\hbar\omega}\ham + \inv{2} \ket{n} =
	\inv{\hbar\omega} \hbar\omega \pab{n + \inv{2}} + \inv{2} = n + 1.
\end{gather*}
Fazę tych stałych możemy wybrać dowolną, a więc też taką, aby współczynniki te były dodatnie:
\begin{align*}
	\vu{a} \ket{n} &= \sqrt{n} \ket{n - 1}, &
	\vu{a}^{\dagger} \ket{n} &= \sqrt{n + 1} \ket{n + 1}.
\end{align*}
Iteratywnie stosując te wzory do stanu podstawowego otrzymujemy wzory ogólne:
\begin{align*}
	\ket{n} &= \inv{\sqrt{n!}} \pab{\vu{a}^{\dagger}}^n \ket{0}, &
	\bra{n} &= \inv{\sqrt{n!}} \bra{0} \pab{\vu{a}}^n.
\end{align*}

Aby znaleźć dokładną formę stanu podstawowego, a następnie wyższych możemy wykorzystać własność tego stanu:
\begin{gather*}
	\vu{a} \ket{0} = 0, \\
	\braket{Q}{0} = f_{0}\pab{Q}, \\
	\pab{Q + \pdv{}{q}} f_{0}\pab{Q} = 0.
\end{gather*}
Rozwiązanie wraz ze stałą normującą ma postać:
\begin{equation}\label{eq:osc_harm_stan_podst}
	f_{0} \pab{Q} = \pi^{-\frac{1}{4}} e^{-Q^2}.
\end{equation}
